\documentclass[11pt,letterpaper]{article}

% Packages
\usepackage[margin=0.6in]{geometry}
\usepackage{enumitem}
\usepackage{hyperref}
\usepackage{titlesec}
\usepackage{xcolor}
\usepackage{fontspec}
\usepackage{tabularx}
\usepackage{ulem}
\usepackage{tcolorbox}
\tcbuselibrary{skins}
\usepackage{microtype}
\usepackage{xeCJK}
\tolerance=1000
\emergencystretch=1em
\setCJKmainfont{Songti SC}[BoldFont=Songti SC Bold]

% Colors
\definecolor{linkblue}{RGB}{0,102,204}

\newtcolorbox{statementbox}{
    enhanced,
    frame hidden,
    colback=linkblue!5,
    borderline west={2pt}{0pt}{linkblue},
    boxrule=0pt,
    left=6pt, right=6pt, top=3pt, bottom=3pt,
    sharp corners
}

% Hyperlink setup
\hypersetup{
    colorlinks=true,
    linkcolor=linkblue,
    urlcolor=linkblue,
    pdfborder={0 0 0}
}

% Section formatting - clean small caps with refined rule
\titleformat{\section}
    {\normalfont\large\scshape\color{linkblue}}
    {}{0em}{}[\vspace{2pt}\titlerule]
\titlespacing*{\section}{0pt}{12pt}{6pt}

% Custom commands
\newcommand{\myname}[1]{\textbf{\uline{#1}}}
\newcommand{\dateright}[1]{\hfill #1}
\newcommand{\institution}[1]{\textcolor{linkblue}{\textbf{#1}}}
\newcommand{\project}[1]{\textbf{\textcolor{black!80}{#1}}}
\newcommand{\advisor}[1]{\textit{Advisor: #1}}

% Remove page numbers
\pagenumbering{gobble}

% Adjust list spacing - standard academic CV indentation
\setlist[itemize]{leftmargin=2em, nosep, topsep=3pt, itemsep=2pt}
\setlist[enumerate]{leftmargin=2em, nosep, topsep=3pt, itemsep=3pt}

\begin{document}

% Updated date in top right corner
\begin{flushright}
\vspace{-1em}
\footnotesize\textit{Updated March 2026}
\end{flushright}

\vspace{-1em}

% Header
\begin{center}
{\LARGE\bfseries Muhua (黄牧华) Huang}\\[0.4em]
{\small\color{gray}
\href{mailto:muhua@stanford.edu}{muhua@stanford.edu}
\quad$\cdot$\quad
\href{https://muhua-h.github.io/}{muhua-h.github.io}}
\end{center}

\vspace{-0.3em}

% Research statement
\begin{statementbox}
\noindent I study how AI systems fundamentally differ from \textbf{human intelligence} and how they, as \textbf{social actors}, reshape \textbf{collaboration} and \textbf{collective decision-making}. My research combines \textbf{computational social science}, \textbf{psychometrics}, and \textbf{generative agent-based modelling (GABM)} to investigate AI's unique cognitive characteristics, uncover new insights about human behavior from AI representations, and understand how AI agents alter \textbf{interaction dynamics} within \textbf{groups and organizations}.
\end{statementbox}

\vspace{0.5em}

% Skills summary
\noindent\textcolor{linkblue}{\textbf{Research:}} Human-AI Collaboration, Organizational Behavior, AI Alignment,
Computational Social Science\\
\noindent\textcolor{linkblue}{\textbf{Engineering:}} GABM, Large-Scale LLM
Experimentation, Fine-tuning, Agent, High Performance Computing\\
\noindent\textcolor{linkblue}{\textbf{Analytical:}} Psychometrics, Embedding Models, Social Network
Analysis, ML/NLP, Behavioral Experiment

%----------------------------------------------------------------------------------------
% EDUCATION
%----------------------------------------------------------------------------------------
\section*{Education}

\noindent
\begin{tabularx}{\textwidth}{Xr}
\institution{Stanford Graduate School of Business} & 2025.9 -- 2030.6 \\
\quad PhD in Organizational Behavior & \\[0.5em]
\institution{The University of Chicago} & 2023.9 -- 2025.6 \\
\quad MA in Computational Social Science; GPA: 3.9/4.0 & \\[0.5em]
\institution{The University of British Columbia (UBC)} & 2018.9 -- 2023.5 \\
\quad BA in Computer Science \& Psychology (Honours); GPA: 4.2/4.3 & \\[0.5em]
\institution{International AI Evaluation Programme Fellow} & 2026.1 -- 2026.5 \\[0.5em]
\institution{Summer Institute in Computational Social Science (SICSS)} & 2024.6 -- 2024.7 \\[0.5em]
\institution{Google Research Mentorship Program Scholar} & 2020.9 -- 2021.1 \\
\end{tabularx}

%----------------------------------------------------------------------------------------
% PUBLICATIONS
%----------------------------------------------------------------------------------------
\section*{Publications}

\begin{enumerate}[label={[\arabic*]}]
    \item \myname{Huang, M.}, Zhang, X., Soto, C., \& Evans, J. (2026). Designing AI-Agents with Personalities: A Psychometric Approach. \textit{Personality Science}. \href{https://doi.org/10.1177/27000710251406471}{DOI}

    \item Zhang, X., \myname{Huang, M.}, Sun, J., \& Savalei, V. (2025). Improving the Measurement of the Big Five via Alternative Formats for the BFI-2. \textit{Journal of Personality Assessment}. \href{https://doi.org/10.1080/00223891.2025.2531187}{DOI}

    \item Yao, J., Yi, X., Duan, S., Wang, J., Bai, Y., \myname{Huang, M.}, \ldots \& Xie, X. (2025). Value compass benchmarks: A comprehensive, generative and self-evolving platform for LLMs' value evaluation. \textit{Proceedings of the 63rd Annual Meeting of the Association for Computational Linguistics}. \href{https://doi.org/10.18653/v1/2025.acl-demo.64}{DOI}

    \item Kim, H., Yi, X., Bak, J., Yao, J., Lian, J., \myname{Huang, M.}, Duan, S., \& Xie, X. (2025). The Road to Artificial SuperIntelligence: A Comprehensive Survey of Superalignment. \textit{SuperIntelligence - Robotics - Safety \& Alignment}, 2(1). \href{https://doi.org/10.70777/si.v2i1.13963}{DOI}

    \item Savalei, V., \& \myname{Huang, M.} (2025). Fit Indices Are Insensitive to Multiple Minor Violations of Perfect Simple Structure in Confirmatory Factor Analysis. \textit{Psychological Methods}. \href{https://dx.doi.org/10.1037/met0000718}{DOI}

    \item Laurin, K., Engstrom, H., \& \myname{Huang, M.} (2024). What will my life be like when I'm 25? How social class predicts kids' answers to this question, and how their answers predict their futures. \textit{Journal of Social Issues}. \href{https://doi.org/10.1111/josi.12650}{DOI}
\end{enumerate}

%----------------------------------------------------------------------------------------
% WORKING PAPERS
%----------------------------------------------------------------------------------------
\section*{Working Papers}

\begin{enumerate}[label={[\arabic*]}]
    \item \myname{Huang, M.}, \& Guilbeault, D. (In Prep). AI as Alien Intelligence: Shared Vocabulary Masks Divergent Category Systems in Human and AI Collectives. \textit{Manuscript available upon request.}

    \item \myname{Huang, M.}, \& Evans, J. (In Press). Institutions as cached computation for resource-rational negotiation. \textit{Behavioral and Brain Sciences}. \href{https://doi.org/10.1017/S0140525X2510174X}{DOI}

    \item Bai, Y., Duan, S., \myname{Huang, M.}, Yao, J., Liu, Z., Zhang, P., \ldots \& Xie, X. (Under Review). IROTE: Human-like Traits Elicitation of Large Language Model via In-Context Self-Reflective Optimization. [\href{https://arxiv.org/abs/2508.08719}{Preprint}]

    \item Zhang, H., \myname{Huang, M.}, \& Wang, J. (Under Review). Evolving Collective Cognition in Human-Agent Hybrid Societies: How Agents Form Stances and Boundaries. [\href{https://arxiv.org/abs/2508.17366}{Preprint}]

    \item Kim, H., Yi, X., Yao, J., \myname{Huang, M.}, Bak, J., Evans, J., \& Xie, X. (2025). Research on Superalignment Should Advance Now with Parallel Optimization of Competence and Conformity. [\href{https://arxiv.org/abs/2503.07660}{Preprint}]

\end{enumerate}

%----------------------------------------------------------------------------------------
% CONFERENCE PRESENTATIONS
%----------------------------------------------------------------------------------------
\section*{Conference Presentations}

\noindent\textsuperscript{$\dagger$} first author / \textsuperscript{*} solo

\vspace{0.3em}
\begin{itemize}[itemsep=4pt]
    \item \textsuperscript{*} \textbf{Designing AI-Agents with Personalities: A Psychometric Approach}\\
    {\scriptsize \textit{APA}, Seattle (Aug 2024, Poster); \textit{CPA}, Montréal (Jun 2026, w/ X. Zhang); \textit{ISDSA}, Beijing (Jul 2026, Speed Talk, w/ X. Zhang)}
    \item \textsuperscript{$\dagger$} \textbf{Beyond Anthropomorphism: Unveiling Unique Value Structure of Large Language Models} (w/ P. Biedma, X. Yi, L. Huang, M. Sun, J. Evans, X. Xie)\\
    {\scriptsize \textit{IC2S2}, Norrköping (Jul 2025)}
    \item \textsuperscript{$\dagger$} \textbf{Re-Discovering the Big Five: Using LLM Embeddings to Understand Personality Structure} (w/ N. Huang)\\
    {\scriptsize \textit{APA}, Denver (Aug 2025)}
    \item \textsuperscript{$\dagger$} \textbf{Improving the Measurement of the Big Five via Alternative Formats for the BFI-2} (w/ X. Zhang, V. Savalei)\\
    {\scriptsize \textit{APA}, Washington DC (Aug 2023, Poster); \textit{IMPS}, Prague (Jul 2025, w/ J. Sun); \textit{CPA}, Montréal (Jun 2026)}
    \item \textsuperscript{$\dagger$} \textbf{Social Class Differences in Children's Hopes and Expectations for the Future} (w/ H. Engstrom, K. Laurin)\\
    {\scriptsize \textit{APS} (May 2021, Poster); \textit{SPSP}, San Francisco (Feb 2022, Talk); \textit{Society in the Classroom}, London (Jul 2022, Talk)}
    \item \textbf{Understanding Silicon Sampling through a Psychometric Lens} (w/ C. F. Falk, X. Zhang, N. Guenole, W. Li)\\
    {\scriptsize \textit{CPA}, Montréal (Jun 2026)}
    \item \textbf{Can SEM Fit Indices Distinguish Between CFA and EFA Data Structures?} (w/ V. Savalei)\\
    {\scriptsize \textit{SMEP}, Iowa City (Oct 2023)}
    \item \textbf{Gauging Student Engagement with an XAI Interface via Eye-tracking} (w/ C. Conati, R. Murali)\\
    {\scriptsize \textit{IJCAI Workshop on XAI} (Jun 2022)}
\end{itemize}

%----------------------------------------------------------------------------------------
% INVITED TALKS
%----------------------------------------------------------------------------------------
\section*{Invited Talks}

\begin{enumerate}[label={[\arabic*]}]
    \item Yang, L., Sun, L., \myname{Huang, M.}, Wang, J., Jiang, R., \& Xu, F. (November 14, 2024). Panel discussion: LLM-driven social science and generative agents. \textit{2024 MSR Asia TAB Workshop: Shaping the Future with Societal AI}. Microsoft Research Asia, Beijing, China.

    \item \myname{Huang, M.} (January 13, 2025). Designing LLM-agents with personalities: A psychometric approach. Invited talk at the \textit{Quantitative Methods Forum}, York University, Toronto, Canada.
\end{enumerate}

%----------------------------------------------------------------------------------------
% INTERNSHIP
%----------------------------------------------------------------------------------------
\section*{Industry Experience}

\noindent\textbf{Research Intern} $|$ Microsoft Research Asia \hfill 2024.7 -- 2024.10\\
\noindent\advisor{Dr. Xiaoyuan Yi, Dr. Xing Xie}\\[0.3em]
%\noindent\project{On the Dynamics of LLMs' Values as A Community}
\begin{itemize}
    \item Built a multi-agent simulation platform where hundreds of LLM agents interact, form social networks, and self-organize governance structures (e.g., constitutions) without external guidance.
    \item Grounded agent behavior in Schwartz's Theory of Basic Values and used Social Network Analysis to trace how value-diverse communities form, evolve, and stabilize over time.
    \item Discovered that moderate value diversity fosters creativity and stability in agent communities, while extreme heterogeneity induces instability --- with implications for designing diverse human-AI teams.
    \item Investigated social emergence (unique language, cultural practices, institutions) in agent communities using GraphRAG and NLP.
    \item Contributed to interdisciplinary projects on AI values and alignment, including LLM value evaluation benchmarks (ACL 2025), psychometrically grounded trait elicitation, and a survey of superalignment.
\end{itemize}

%----------------------------------------------------------------------------------------
% RESEARCH EXPERIENCE
%----------------------------------------------------------------------------------------
\section*{Research Experience}

\noindent\textbf{PhD Student Researcher} $|$ Stanford Graduate School of Business \hfill 2025.10 -- Present\\
\noindent\advisor{Dr. Douglas Guilbeault}\\[0.3em]
\noindent\project{AI as Alien Intelligence: Shared Vocabulary Masks Divergent Category Systems in Human and AI Collectives}
\begin{itemize}
    \item Engineered a large-scale multi-agent coordination experiment ($\sim$375,000 AI trials with GPT-4o agents in network structures of $N$=2 to 50), replicating a human category-formation paradigm.
    \item Discovered ``representational relativity'': despite comparable behavioral performance, human and AI collectives resolve ambiguity in systematically different places ($r = -0.130$, $p < 0.001$).
    \item Showed AI labels are disproportionately visual-dominant and concrete, while human labels draw on richer multimodal experience --- divergence driven by grounding modality, not training data.
    \item Found that AI agents strategically invent novel vocabulary to express distinctions human language cannot; forcing AI to use only human labels reduces performance by 12.7\%.
\end{itemize}

\vspace{0.5em}

\noindent\textbf{Master's Thesis} $|$ University of Chicago \hfill 2023.9 -- 2024.6\\
\noindent\advisor{Dr. James Evans}\\[0.3em]
\noindent\project{Designing LLM-Agents with Personalities: A Psychometric Approach}
\begin{itemize}
    \item Developed the first psychometrically validated framework for assigning Big Five personalities to LLM agents, with controllable, fine-grained design across multiple models (GPT-3.5, GPT-4, Claude 3).
    \item Validated that personality-assigned agents exhibit predictable behavioral differences in moral reasoning, risk-taking, and cooperation (personality explains 38--52\% of variance in agent decisions).
    \item Collected data from 350 human participants for direct human--agent comparison, demonstrating potential for using personality-designed agents in behavioral research.
    \item Published in \textit{Personality Science}; methodology enables designing AI agents for specific collaborative roles.
\end{itemize}

\vspace{0.5em}

\noindent\textbf{Research Assistant, Quantitative Psychology} $|$ University of British Columbia \hfill 2020.4 -- 2023.8\\
\noindent\advisor{Dr. Victoria Savalei}\\[0.3em]
\noindent\project{Exploring the Factor Structure of the BFI-2 in Alternative Scale Formats}
\begin{itemize}
    \item Compared four personality scale formats using SEM; found that alternative formats (Expanded, Item-Specific) significantly improve reliability and model fit over Likert. Published in \textit{Journal of Personality Assessment}.
    \item Honours thesis nominated for multiple Canadian national undergraduate research awards (Belkin, CPA).
\end{itemize}
\noindent\project{SEM Fit Indices' Sensitivity to Omitted Crossloadings in CFA}
\begin{itemize}
    \item Built an R Shiny simulation app to examine fit index performance when CFA models are applied to data with EFA-like structure; contributed to manuscript published in \textit{Psychological Methods}.
\end{itemize}

\vspace{0.5em}

\noindent\textbf{Research Assistant, Social Psychology} $|$ University of British Columbia \hfill 2020.8 -- 2021.9\\
\noindent\advisor{Dr. Kristin Laurin}\\[0.3em]
\noindent\project{Socioeconomics and Imagined Future Study}
\begin{itemize}
    \item Applied ML (Random Forest, Neural Network) and NLP (Topic Modelling, LSTM, Word2Vec) to analyze 10,000+ youth essays, revealing how socioeconomic background shapes imagined futures. Published in \textit{Journal of Social Issues}.
    \item Supervised 20 research assistants and built data processing protocols including a word distribution and performance tracking system.
\end{itemize}

\vspace{0.5em}

\noindent\textbf{Research Assistant, Human-Centered AI} $|$ University of British Columbia \hfill 2021.10 -- 2022.9\\
\noindent\advisor{Dr. Christina Conati}\\[0.3em]
\noindent\project{Personalized Explainable AI for Intelligent Tutoring System Study}
\begin{itemize}
    \item Evaluated how personalizing AI explanations to individual cognitive profiles (perceptual speed, personality) affects learning outcomes, using eye-tracking and linear mixed models.
    \item Found that tailoring AI tutoring interfaces to users' cognitive abilities improves engagement and effectiveness --- demonstrating the value of adaptive human-AI interaction design.
\end{itemize}

%----------------------------------------------------------------------------------------
% WORK EXPERIENCE
%----------------------------------------------------------------------------------------
\section*{Professional Experience}

\noindent\textbf{Computational Social Science Workshop Coordinator} $|$ University of Chicago \hfill 2024.10 -- 2025.6
\begin{itemize}
    \item Organized the weekly MACSS Computational Social Science Workshop, featuring speakers on advanced computational methods and interdisciplinary research.
\end{itemize}

\vspace{0.5em}

\noindent\textbf{Deep Learning Textbook Editor} $|$ University of Chicago \hfill 2023.10 -- 2024.10
\begin{itemize}
    \item Proofread and edited ``Thinking with Deep Learning'' by Dr. James Evans; designed LaTeX templates for the textbook and for Journal of Social Computing.
\end{itemize}

\vspace{0.5em}

\noindent\textbf{Teaching Assistant} $|$ University of British Columbia \hfill 2021.01 -- 2022.12
\begin{itemize}
    \item TA for Human Computer Interaction Methods and Systematic Program Design; rated 4.7/5 by students and received three return offers.
\end{itemize}

\vspace{0.5em}

\noindent\textbf{Crisis Respondent} $|$ Kids Help Phone \hfill 2019.09 -- 2020.08
\begin{itemize}
    \item Provided over 230 hours of crisis counselling to 550+ individuals in severe mental distress.
\end{itemize}

%----------------------------------------------------------------------------------------
% AWARDS & HONORS
%----------------------------------------------------------------------------------------
\section*{Honors \& Awards}

\noindent
\begin{tabularx}{\textwidth}{Xr}
Stanford Graduate School of Business Doctoral Fellowship & 2025 -- 2030 \\
Stanford University EDGE Doctoral Fellowship (\$12,800) & 2025 \\
UChicago Outstanding Thesis Award (\$1,000) & 2025 \\
Microsoft Research Asia Outstanding Intern Award (Stars of Tomorrow Certificate) & 2024 \\
OpenAI Researcher Access Program (\$3,500) & 2024 \\
SICSS-Beijing Merit Based Scholarship (¥3,600) & 2024 \\
UChicago Computational Social Science Research Poster Competition (Honorable Mention, 2\textsuperscript{nd} Place) & 2024 \\
UChicago Quadrangle Research Scholarship (\$80,000) & 2023, 2024 \\
UChicago Social Sciences Promise \& Merit Scholarship (\$10,000) & 2023, 2024 \\
Quinn Research Assistantship Award (\$8,300) & 2021 \\
UBC International Community Achievement Award (\$5,000) & 2021 \\
UBC Trek Excellence Scholarship (\$2,000) & 2019, 2020 \\
UBC Faculty of Arts International Student Scholarship (\$8,800) & 2019, 2020, 2022 \\
\end{tabularx}

\end{document}
